\begin{lemma}
For a fixed set \(I\) with
\[
|I|=k=\noderef{TwoBiteNaturalIndependenceScale}(1+\varepsilon,n),
\]
assume
\(0\le\varepsilon_1\le1\),
\(\noderef{ParameterHierarchyEventualComparisons}
(\varepsilon,\varepsilon_1,\varepsilon_2,n_0)\), \(n>n_0\), and the rounded
sample-parameter identities
\[
m=\noderef{TwoBiteNaturalReducedVertexCount}(n),
\qquad
p=\noderef{TwoBiteEdgeProbability}\!\left(\frac12,n\right).
\]
Then the conditional probability
\(\noderef{TwoBiteConditionalProbability}(n,m,p;-|-)\) with numerator
\[
I\text{ independent in }\noderef{TwoBiteFinalGraph}(\omega)_0
\quad\text{and}\quad
\noderef{TwoBiteRegularityEvent}(k,\varepsilon,\varepsilon_1,\varepsilon_2,
\omega)
\]
and denominator
\[
\noderef{TwoBiteProjectionSizeEvent}(I,k,\ell_R,\ell_B,\omega)
\]
is at most
\[
\exp\left(-p\left(f(\ell_R,\ell_B)-\varepsilon^3k^2\right)\right).
\]
Here
\[
f(\ell_R,\ell_B)=
\noderef{ProjectionOpenPairFunction}(\ell_R,\ell_B,k,p,n).
\]
\end{lemma}
\begin{proof}
Let
\[
S(\omega)=
\noderef{TwoBiteProjectionSizeEvent}(I,k,\ell_R,\ell_B,\omega)
\]
be the projection-size event in the denominator, and let
\[
R(\omega)=
\noderef{TwoBiteRegularityEvent}(k,\varepsilon,\varepsilon_1,\varepsilon_2,\omega)
\]
be the regularity event carried in the numerator.  Thus the probability to be
bounded is
\(\Pr(\mathcal B_I\cap R\mid S)\), where \(\mathcal B_I\) is final-graph
independence of \(I\).

We first verify the pointwise hypothesis required by
\(\noderef{FixedSetExposureExceptionalChoices}\).  Fix an arbitrary sample
\(\omega\), and assume
\[
H_{\mathcal R}:R(\omega),\qquad H_{\mathcal S}:S(\omega).
\tag{1}
\]
Apply \(\noderef{FixedSetOpenPairsAfterConditioning}\) to this same
\(\omega\) and this same set \(I\), with the same integers
\(n,m,k,\ell_R,\ell_B\), the same real parameters
\(p,\varepsilon,\varepsilon_1,\varepsilon_2\), and the same \(n_0\).
The hypotheses passed to it are exactly
\[
0\le\varepsilon_1,\quad
\noderef{ParameterHierarchyEventualComparisons}
(\varepsilon,\varepsilon_1,\varepsilon_2,n_0),\quad
n>n_0,
\]
\[
m=\noderef{TwoBiteNaturalReducedVertexCount}(n),\quad
p=\noderef{TwoBiteEdgeProbability}\!\left(\frac12,n\right),\quad
|I|=k=\noderef{TwoBiteNaturalIndependenceScale}(1+\varepsilon,n),
\]
together with \(H_{\mathcal R}\) and \(H_{\mathcal S}\) from (1).  Therefore
this arbitrary sample in \(R\cap S\) satisfies
\[
\noderef{ProjectionOpenPairFunction}(\ell_R,\ell_B,k,p,n)
-10\varepsilon_1k^2
\le
\left|\noderef{TwoBiteStagedOpenPairs}(\omega,\varepsilon,I)\right|.
\tag{2}
\]
Because \(\omega\), \(H_{\mathcal R}\), and \(H_{\mathcal S}\) were arbitrary,
(2) is precisely the pointwise implication
\[
\forall\omega,\;
R(\omega)\to S(\omega)\to
\noderef{ProjectionOpenPairFunction}(\ell_R,\ell_B,k,p,n)
-10\varepsilon_1k^2
\le
\left|\noderef{TwoBiteStagedOpenPairs}(\omega,\varepsilon,I)\right|.
\tag{3}
\]

Now apply \(\noderef{FixedSetExposureExceptionalChoices}\) with the same
fixed set \(I\), the same integers \(n,m,k,\ell_R,\ell_B\), the same real
parameters \(p,\varepsilon,\varepsilon_1,\varepsilon_2\), and the same
\(n_0\).  Its scalar and hierarchy hypotheses are exactly the present
hypotheses:
\[
0\le\varepsilon_1,\quad \varepsilon_1\le1,\quad
\noderef{ParameterHierarchyEventualComparisons}
(\varepsilon,\varepsilon_1,\varepsilon_2,n_0),\quad n>n_0,
\]
\[
m=\noderef{TwoBiteNaturalReducedVertexCount}(n),\qquad
p=\noderef{TwoBiteEdgeProbability}\!\left(\frac12,n\right),
\]
\[
|I|=k,\qquad
k=\noderef{TwoBiteNaturalIndependenceScale}(1+\varepsilon,n).
\]
Its remaining hypothesis is exactly the pointwise implication (3).  The
conditional probability in its conclusion has denominator \(S\) and numerator
event \(\mathcal B_I\cap R\), exactly as here.  The child theorem therefore
gives directly
\[
\Pr(I\text{ is independent in }\noderef{TwoBiteFinalGraph}
     \text{ and }R\mid S)
\le
\exp\left(-p\left(
\noderef{ProjectionOpenPairFunction}(\ell_R,\ell_B,k,p,n)
-\varepsilon^3k^2\right)\right),
\]
which is exactly the asserted bound.  No further exponent estimate or
probability counting is introduced at this aggregation step.
\end{proof}
